%%=============================================================================
%% Samenvatting
%%=============================================================================

% De "abstract" of samenvatting is een kernachtige (~ 1 blz. voor een
% thesis) synthese van het document.
%
% Een goede abstract biedt een kernachtig antwoord op volgende vragen:
%
% 1. Waarover gaat de bachelorproef?
% 2. Waarom heb je er over geschreven?
% 3. Hoe heb je het onderzoek uitgevoerd?
% 4. Wat waren de resultaten? Wat blijkt uit je onderzoek?
% 5. Wat betekenen je resultaten? Wat is de relevantie voor het werkveld?
%
% Daarom bestaat een abstract uit volgende componenten:
%
% - inleiding + kaderen thema
% - probleemstelling
% - (centrale) onderzoeksvraag
% - onderzoeksdoelstelling
% - methodologie
% - resultaten (beperk tot de belangrijkste, relevant voor de onderzoeksvraag)
% - conclusies, aanbevelingen, beperkingen
%
% LET OP! Een samenvatting is GEEN voorwoord!

%%---------- Nederlandse samenvatting -----------------------------------------
%
% Als je je bachelorproef in het Engels schrijft, moet je eerst een
% Nederlandse samenvatting invoegen. Haal daarvoor onderstaande code uit
% commentaar.
% Wie zijn bachelorproef in het Nederlands schrijft, kan dit negeren, de inhoud
% wordt niet in het document ingevoegd.

%\IfLanguageName{english}{%
%\selectlanguage{dutch}
%\chapter*{Samenvatting}
%\lipsum[1-4]
%\selectlanguage{english}
%}{}

%%---------- Samenvatting -----------------------------------------------------
% De samenvatting in de hoofdtaal van het document

%TODO: aanvullen en nakijken
\chapter*{\IfLanguageName{dutch}{Samenvatting}{Abstract}}

% TODO: pas de concrete cijfers en modelresultaten aan 
% zodra de definitieve evaluatie beschikbaar is.
POI-databanken (Points-of-Interest) vormen de ruggengraat van 
locatiegebaseerde analyses zoals marktstudies, bezoekersstroomonderzoek 
en concurrentie-inzicht. Een veelvoorkomend probleem is echter dat de 
operationele status van POI's — of een locatie open dan wel gesloten is — 
frequent verouderd of onvolledig is in deze databanken. Dit leidt tot 
foutieve beslissingen en ondermijnt de betrouwbaarheid van 
locatiegebaseerde toepassingen. Huidige monitoringmethoden zijn 
grotendeels handmatig, wat resulteert in een aanzienlijke vertraging 
tussen de werkelijke toestand van een POI en de registratie ervan.

Deze bachelorproef onderzoekt hoe artificiële intelligentie kan worden 
ingezet om automatisch de (historische) open of gesloten status van 
POI's te detecteren, met supermarkten als specifiek toepassingsdomein. 
Hiervoor werd een proof-of-concept ontwikkeld op basis van de interne 
POI-databank van Accurat, aangevuld met historische statusregistraties 
en gestructureerde openingstijden voor België, Nederland, Luxemburg 
en Frankrijk.

De methodologie omvat een uitgebreide feature engineering fase waarin 
temporele, structurele en historisch gedragsmatige kenmerken werden 
geconstrueerd, waaronder de \texttt{inactiviteit\_ratio} en de 
\texttt{poi\_leeftijd\_maanden}. Drie machine learning modellen werden 
getraind en met elkaar vergeleken: Logistische Regressie als 
referentiemodel, Random Forest en XGBoost. De modellen werden 
geëvalueerd op basis van accuracy, precision, recall en F1-score, 
waarbij de detectie van gesloten vestigingen als primaire 
evaluatiedoelstelling werd gehanteerd.

% TODO: vervang onderstaande zin door de definitieve bevinding 
% op basis van de modelresultaten.
De resultaten tonen aan dat [TODO: bv. \textit{ensemble-methoden 
    zoals XGBoost het sterkst presteren met een F1-score van X\% voor 
    de gesloten klasse}], wat bevestigt dat supervised classificatie 
een praktisch inzetbare aanpak vormt voor POI-statusdetectie op 
basis van intern beschikbare structurele data. Tegelijkertijd 
worden de grenzen van deze aanpak blootgelegd: het ontbreken van 
dynamische mobiliteitsbronnen, geospatiale bias ten gevolge van de 
onevenwichtige geografische dataverdeling en het cold-start probleem 
voor jonge vestigingen blijven structurele uitdagingen.

Dit onderzoek levert een werkend prototype op, aangevuld met een 
vergelijkende evaluatie van de gebruikte modellen en concrete 
aanbevelingen voor vervolgonderzoek. Voor organisaties zoals Accurat 
biedt het een waardevol kader voor het automatisch up-to-date houden 
van POI-data, wat de kwaliteit van locatiegebaseerde analyses 
ten goede komt.
